% !TEX program = xelatex
\documentclass[UTF8,a4paper,zihao=-4,oneside,openany,fontset=fandol]{ctexbook}
\usepackage[margin=23mm,top=23mm,bottom=24mm,headheight=16pt]{geometry}
\usepackage{fontspec,graphicx,booktabs,longtable,tabularx,array,enumitem,fancyhdr,xcolor,fvextra,needspace,float}
\definecolor{InkBlue}{HTML}{234B69}
\definecolor{InkPale}{HTML}{F0F5F8}
\definecolor{InkGray}{HTML}{53606D}
\usepackage[unicode,bookmarksnumbered,colorlinks,linkcolor=InkBlue,urlcolor=InkBlue]{hyperref}
\setmainfont{texgyrepagella-regular.otf}[BoldFont=texgyrepagella-bold.otf,ItalicFont=texgyrepagella-italic.otf,BoldItalicFont=texgyrepagella-bolditalic.otf]
\setsansfont{texgyreheros-regular.otf}[BoldFont=texgyreheros-bold.otf]
\setmonofont{lmmono10-regular.otf}[Scale=MatchLowercase]
\ctexset{chapter={format=\Huge\sffamily\bfseries\color{InkBlue},beforeskip=8pt,afterskip=22pt},section={format=\Large\sffamily\bfseries\color{InkBlue}},subsection={format=\large\sffamily\bfseries}}
\setlength{\parskip}{4pt}
\setlength{\emergencystretch}{2em}
\linespread{1.15}
\setlist{nosep,topsep=6pt,leftmargin=2em,itemsep=3pt}
\setcounter{tocdepth}{1}
\setcounter{secnumdepth}{1}
\raggedbottom
\pagestyle{fancy}
\fancyhf{}
\fancyhead[L]{\small\sffamily InkStruct 用户手册}
\fancyhead[R]{\small\nouppercase{\leftmark}}
\fancyfoot[L]{\footnotesize\color{InkGray}采集、生成与排版篇 · 2026-11-04}
\fancyfoot[R]{\thepage}
\renewcommand{\headrulewidth}{0.3pt}
\fancypagestyle{plain}{\fancyhf{}\fancyfoot[L]{\footnotesize InkStruct 用户手册}\fancyfoot[R]{\thepage}\renewcommand{\headrulewidth}{0pt}}
\newcommand{\ui}[1]{\textbf{「#1」}}
\newcommand{\code}[1]{\allowbreak{\small\ttfamily\detokenize{#1}}\allowbreak}
\newcommand{\entry}[1]{\noindent\textcolor{InkBlue}{\textbf{操作入口：}}#1\par}
\newcommand{\note}[1]{\par\noindent\colorbox{InkPale}{\parbox{\dimexpr\linewidth-2\fboxsep}{\small #1}}\par}
\newcommand{\shot}[3]{\begin{figure}[H]\centering\includegraphics[width=#2\linewidth]{assets/#1}\caption{#3}\end{figure}}
\DefineVerbatimEnvironment{Source}{Verbatim}{fontsize=\small,samepage=true,breaklines=true,breakanywhere=true,frame=single,rulecolor=\color{InkBlue},framesep=3mm,xleftmargin=0pt}
\newcolumntype{Y}{>{\raggedright\arraybackslash}X}
\newcolumntype{P}[1]{>{\raggedright\arraybackslash}p{#1}}
\hypersetup{pdftitle={InkStruct 用户手册：采集、生成与排版篇},pdfauthor={InkStruct},pdfsubject={项目、Markdown 混排、模板设计与写字机导出}}
\begin{document}
\begin{titlepage}
\thispagestyle{empty}
\vspace*{18mm}
{\sffamily\color{InkBlue}\fontsize{42}{48}\selectfont InkStruct\par}
\vspace{8mm}
{\Huge\bfseries 用户手册\par}
\vspace{5mm}
{\Large 采集、生成与排版篇\par}
\vspace{10mm}
\noindent\textcolor{InkBlue}{\rule{\textwidth}{1.2pt}}
\vspace{8mm}
{\large 从 Markdown 源码到手写页面与写字机 SVG\par}
\vspace{7mm}
\noindent 本册介绍统一应用的采集、字形生成、字体发布和编辑排版流程。\par
\vfill
\noindent\begin{tabular}{@{}ll@{}}
手册版本 & 0.1.120 本机测试版\\
适用界面 & Windows 统一应用 0.1.120 本机测试版\\
阅读对象 & 文档编写者、模板设计者与协作 Agent\\
当前详述 & 采集、生成、字体、编辑排版与导出\\
后续扩展 & 设备验收、局部重生成与字体校准\\
\end{tabular}
\vspace{14mm}
\end{titlepage}
\frontmatter
\chapter*{如何阅读本手册}
\addcontentsline{toc}{chapter}{如何阅读本手册}
InkStruct 编辑器将 Markdown 源码、实际手写轨迹和毫米页面结合起来。左侧写内容，右侧检查排版；正式导出前，程序检查缺字、尺寸与可书写内容。用系统字体显示在输入框中的文字，不等于已经拥有它的手写字形。

本手册依据当前编辑器界面和实现编写。图示采用本项目的实际测试界面；示例数字仅用于学习操作，不代表实验测量结果。手册中的普通印刷字体用于讲解，与应用导出的手写资源不同。

\section*{按任务选择阅读路径}
\par\noindent\begin{tabularx}{\linewidth}{@{}P{36mm}Y@{}}
\toprule
我要完成的任务 & 建议阅读\\\midrule
第一次打开应用 & 第\ref{ch:basics}章认识文件，第\ref{ch:quick}章完成数字示例。\\
排版一份报告 & 第\ref{ch:workspace}至\ref{ch:math}章，按需要加入表格与公式。\\
在预印表单上填写 & 第\ref{ch:templates}章，先建模板，再绑定内容。\\
检查生成字形好坏 & 第\ref{ch:resources}章，导入资源并使用字形评价。\\
交给写字机 & 第\ref{ch:export}章，区分参考页面与机器文件。\\
让 Agent 帮忙设计模板 & 附录\ref{ch:agent}，交付模板、项目与适配报告。\\
\bottomrule
\end{tabularx}

\section*{统一手册的组织方式}
全书按四个分部组织。第一部放置所有模块共享的项目与资源概念；第二部完整介绍当前编辑排版；第三、四部分别介绍生成和采集的操作与交接。后续扩展在各分部中增加操作章节，不将生成参数、设备校准等内容塞入编辑器章节。

\note{统一应用通过左侧导航提供采集、生成和排版，底部依次为样本库、字体库、模板库和设置。生成默认 v6.2 exp，v6.1 exp 和 v4 adaptive 可选；具体步骤见第\ref{ch:generation}章。}

\section*{约定}
\ui{方括号式文字}表示界面按钮、菜单或选项；等宽框展示可输入的源码。尺寸默认使用毫米。macOS 快捷键以 Command 表示，Windows 以 Ctrl 表示；有差异时以应用菜单显示为准。窗口图示来自 macOS/Qt，Windows 的原生对话框外观可能不同。

本册源码适用应用 0.1.88，部分历史图示沿用 macOS 外观。Windows 启动时使用独立工作区，在界面重新选择本机样本目录。Windows 的实际验收结果以项目内的 Windows 初始化记录为准。真实数位板、驱动、中文输入法组合输入和写字机纸面效果仍需实物验证。

\paragraph{历史排版更新（0.1.43）}本次主应用合入两端对齐、中文自然书写v2与排版问题面板。自然书写默认关闭，旧v1效果保留；显式强制SVG只输出已有有效笔迹并附缺口报告，不代表文稿完整通过常规导出检查。应用正文参数时保留本文模板和纸张，明确修改对应项才替换。历史独立0.1.43预览ZIP不包含打包后新增的自然书写v2与问题面板，应以本次主应用包及交付记录为准。

\paragraph{大字体与成品数量（0.1.46）}字体展开总量上限为4 GiB，按成员未压缩大小计算，不是压缩文件体积；单成员、点数和完整校验仍保留。v6.2生成页的\ui{每字成品上限}可选1--64份，默认4份，包含原主字形；真实供体按完整部件均衡尝试，同一供体字可使用不同录制样本，也可有限尝试两个部件同时变化，逐份检查最终整字。来源不足时不凑数。旧字体、项目和冻结选样不自动改写，方法名称不变。详见第\ref{ch:generation}章与第\ref{ch:resources}章。

\paragraph{生成进度（0.1.57）}浅绿区域的宽度表示已处理比例，深绿区域显示各阶段的处理速度。当前速度以字/秒显示，读数位于图内右侧并随横线升降。图下可查看剩余字数、失败数及预计剩余时间。

\paragraph{长文稿预览（0.1.12）}字体候选和字形边界可在内容完全一致时复用；自动恢复副本在后台保存，期间可继续输入。正式保存仍检查文件是否被其他程序修改。缺字和未支持公式仍会阻止成品导出，性能优化不改变此限制。

\paragraph{连续页面预览（0.1.13）}右侧页面纵向连续排列，可直接滚轮浏览；页码随当前位置更新，输入页码可跳转。缩放时保持当前阅读位置，点击其他页的字形仍可定位源码。屏幕图像按显示缩放缓存，矢量导出继续使用原始笔迹。首次进入未缓存页面时会短暂显示绘制提示；快速滚动仅优先绘制当前位置，内存缓存有上限。

\paragraph{统一设置（现行位置0.1.67）}左下角设置是采集与排版设置的统一入口，支持 Ctrl+逗号。应用设置的主题提供跟随系统、浅色、深色，保存后立即生效。公式设置的两个对照预览采用当前主题配色。采集显示只影响屏幕。汉字与标点扩展的笔迹粗细在采集页分别调整。排版的 Ctrl+N 从全局默认新建文档，导入 Markdown 时自动嵌入源码目录内受支持的本地图片；排版设置保存在当前项目。模板和字体默认值分别在模板库和字体库管理。设置左侧按正文、公式、表格、图片、代码和脚注分类；修改先保留为草稿；存为默认只影响新文档并关闭窗口，应用并保存同时更新当前文档且保持窗口打开，取消放弃尚未保存的修改。表格和图片可沿用文档默认或独立设置，修改默认不覆盖独立设置。问题诊断与字形评价仍在排版工作区。

\paragraph{固定浅色界面（0.1.15）}当前版本固定使用浅色界面，不跟随系统深色模式。此设置仅作用于应用显示，不更改操作系统主题，也不改变文档、笔迹或字体参数。深色模式兼容留待后续版本。

\paragraph{全局默认与元素属性（0.1.16）}齿轮保存本机全局默认，供后续新文档使用；已有项目保留原参数。可明确勾选同时用于当前文档。排版右侧的“元素属性”用于选择表格或图片、设置独立参数或恢复沿用本文默认，点击应用后随项目保存。修改全局默认不清除已有独立属性。公式设置窗口按内容最小宽度限制，支持纵向滚动。

\paragraph{字体导入进度（0.1.17）}导入时先复制到临时区域，完整验证后加入字体库，显示当前阶段与已验证候选数。完整字体仍需校验每个成员、轨迹、动态信息、来源和度量；已有同包编号但内容不同的文件会拒绝导入并保留原文件。

\paragraph{并行验证与防误调（0.1.18）}完整字体导入最多使用四个进程分批校验，小字体和低核心环境串行处理，完整检查不变。参数框、参数滑块和下拉框不使用滚轮改值，经过控件时仍可滚动页面；点击和键盘调整保留。正文设置全部展开，公式页移除旧容器箭头，实际字形预览默认显示辅助线。以上界面行为不改写旧文档参数。

\paragraph{字符校正（0.1.20）}表格、图片设置直接陈列全部选项。点击排版预览字符，在右侧调整位移和宽高比例。范围为“仅此处”“本文同一样本”“本文同一字符的全部样本”：例如选中 A2，依次只改此处、全部 A2 或同类 A1/A2 等所有样本。正斜体默认分开，跨类别需明确勾选；界面显示样本及影响次数。批量操作统一更新范围内的校正，旧单字校正不会遮住结果，可整体撤销、重做或恢复。后两种范围可勾选“同时保存到字体包的新版本”，完整校验并保留原包和来源；成功后应用本文校正，当前文档使用原字体快照，不重复叠加，新包可导入后用于后续文档。取消或失败不修改本文，处理期间文档变化不会套用旧操作。校正随项目保存并用于 SVG，字距、换行、原资源和超界检查保留；旧文档按原效果重开。

\paragraph{模板连续正文（0.1.21）}简单模板重复单页正文，自定义模板使用有序起始页组与循环页组。一个 Markdown 同时存放命名字段与唯一正文，正文按编号跨框流动，字段独立固定区。模板制作仍从齿轮进入，本份文档从右侧“模板内容”套用和填写。旧模板及项目保持原版读取。

\paragraph{模板界面整理（0.1.22）}齿轮中的模板管理与排版默认分开。工作台统一草稿编辑，纸张、框与字段使用明确属性；右侧“模板”提供字段表单和适配问题。表单仍修改唯一 Markdown，保留原协议和冻结选样。

\paragraph{英文字距与邻字参考（0.1.23）}正文与段落设置增加按笔迹宽度排列英文字母，排除单字采集的横向位置留白，不缩放笔迹。旧文档和已有默认保持旧字距；新安装使用推荐值。明确勾选将默认排版参数应用到当前文档时，会同时应用设置中的其他排版默认，应先核对。字符校正预览显示同一行真实邻字及字距，淡蓝底标识选中位置，灰色为原字形，蓝色为校正草稿；批量范围同步预示可见邻字变化，应用前不改变文档与冻结选样。

\paragraph{图片高度与汉字公式（0.1.24）}图片默认高 60 mm、宽度按比例，默认页不再提供宽度输入。单图可独立选择尺寸依据。旧文档宽度与个性化属性继续保留，转换为新默认需明确操作。公式直接支持汉字及汉字上下标，例如 \code{U_{模块}}，仍检查真实字体、源码与导出能力。

\paragraph{跳过参考图片导出（0.1.25）}含非 SVG 图片时先提示再导出；确认后保留占位与图题，输出清单记录跳过项。图片资源与文档保留，缺字、公式及越界检查继续生效。

\paragraph{缺失位图引用的导出（0.1.26）}尚未载入的本地 PNG/JPEG 可与已载入位图一起明确确认跳过。弹窗分别计数并可查看文件名；缺失错误保留，其他阻断问题仍需修正。导出保持当前占位、图题、页序、坐标和冻结选样。

\paragraph{模板预览停止（0.1.28）}某对象无法放入整个循环页组时，保留此前已排正文及已引用脚注，明确提示此处及后续正文尚未排版。通过查看问题定位原稿，调整正文框或手动拆分公式；不自动缩放、截断或允许成品导出。较大字体资源目录采用独立有界容量，校验与原文档保留。

\paragraph{标点间距与路径优化（0.1.29）}新文档的中英文标点均按实际笔迹宽度紧凑排列，不保留半字或整字空位。保留最小笔迹间隙、大小、基线与选样；公式及代码沿用原规则。齿轮中的正文与段落提供标点间距选择，应用到旧文档须明确勾选。路径优化默认开启，旧文档保留原值；优化只调整允许的轨迹顺序与方向，不改变字形和坐标。

\paragraph{紧凑公式（0.1.30）}新文档公式消除采集位置带来的横向留白，按真实笔迹边界布置分式、根式和括号，降低默认运算符及结构净距，保留笔宽保护和显式空格。根号和伸缩括号默认使用录制字形；根号钩部与上横线分段适配，括号按样本比例和类型限制宽度。旧文档保留原样式，明确应用新默认后才重排；原字体和选样不改写。

\paragraph{预览与 SVG 导出（0.1.33）}左下角设置的“预览绘制”可选择 CPU（默认）或 GPU，保存后下次启动生效；GPU 失败自动回退 CPU。两者均复用有限 SVG 解析缓存，并避免反复取消相邻页预取。只影响显示，不改变排版和选样。
“导出 SVG”打开选项弹窗：页码留空表示全部，也可输入 1-3,5；可取消表格内内容，仅保留表格线和表外内容。路径优化默认开启。先严格校验原始排版，再以左上角为原点同步缩放输出点坐标、SVG 画布宽高与 viewBox，笔宽不变；输出尺寸随比例改变。原文档纸张、预览和冻结选样不变。缩放后的路径及画布仍须通过校验。定位点为各内容文本框左上角、右下角边界上的 0.1 mm 短线，笔宽 0.01 mm；校正比例也作用于短线。保持冻结快照及原文档页码，取消不输出。
制作字体时，来源齐全且目标未录制时可合成摄氏度（中文句号缩小抬高，加正体 C），短横线复用减号；全角波浪号从两笔约等号各提取一笔，每份约等号提供两份候选。记录来源与变换，不伪造独立采集，不覆盖旧字体和冻结文档。

\paragraph{窄数字占位（0.1.32）}紧凑公式对数字保留最小占位，避免接近竖线的数字1与邻字过密。参照表格默认数字参考高度与占位下限，最小占位为当前层公式字号的0.315倍；只对不足部分居中补白，不缩放笔迹。分式和上下标按各自字号计算，非数字和表格普通数字规则保持。旧版采集留白公式不受影响，已有紧凑公式重新排版时直接使用修复规则，原稿、样式和选样保留。

\paragraph{v6.2 exp 验证范围}新增止部关键笔和竖弯钩语义准入，保留 v6.1 与 v4。user\_001 验证合并字体仅重生孔、正，其余字形保持；孔末笔保留明确参考缺口，不代表全作者字体。已有文档继续使用已冻结资源。


\paragraph{录制符号与翻页（0.1.35）}正文星号 U+002A 与数学星号 U+2217 保留独立字符身份，使用减号录制笔迹缩短后按 0、60、120 度组合；优先三份不同样本，不足时明确记录重复使用。确定性循环配对，每个目标至多与合格减号候选数量相同，不枚举所有组合。公式内星号的原解析含义不变。
数学比号 U+2236 缺失时优先复用全角冒号，再使用半角冒号；数学圆圈 U+2218 缺失时复用中文句号，按数学轴放置，圆圈命令的含义和间距类别不变。整份笔迹等比缩放，不描绘理想点或圆。独立目标优先，记录来源、仿射变换、构造身份和时间平移；不是独立录制。旧字体扩展另存新包，不改旧文档与冻结选样。
排版页页码右侧向上按钮跳至上一页，向下按钮跳至下一页；键盘方向一致，可直接输入页码，连续滚动仍同步当前页。

\paragraph{侧栏与预览（0.1.57）}系统标题栏显示 InkStruct 与应用版本，不再占用独立品牌横栏。左下设置位于模板库下方；Ctrl+逗号保留。Alt+1/2/3/4/5切换采集、生成、排版、字体库、模板库；F1或侧栏右键打开关于。排版预览支持 Ctrl＋滚轮缩放。启动时可直接显示空白 PDF 模板的已校验底图，默认字体在后台校验。依赖字体的页码将在校验后绘出。模板工作台可追加或替换目标页组的 PDF 页面版式。字体库预设直接展示汉字或公式笔迹，点击\ui{编辑试写…}输入自定义内容，切换预设后草稿保留。来源、缺字和完整校验状态始终如实显示。

在两台电脑之间录制样本时，可通过「采集 › 书写者资料」导出和导入同一稳定书写者 ID。返回原设备后，将 JSON、ZIP 或样本目录导入已有书写者。旧汉字批次与另一台电脑录制的标点可在制作字体时核对双方 ID 和样本范围，确认后在本机建立批次级关联。

\paragraph{汉字间距（0.1.85）}在「设置 › 正文与段落」可选整字宽或按笔迹紧凑排列。保存项目后保留所选模式，旧项目按整字宽打开。自然书写的字距变化可独立调节。

\paragraph{样本库（0.1.86）}在「采集」打开样本库，按书写者、字符、类别或任务筛选，预览原始笔迹并选择导出。接收端选择目标书写者，核对离线采集包中的新增、重复和冲突后导入。

\paragraph{段落对齐区块（0.1.86）}分别在独立源码行使用 \verb|\begin{paragraphalign}{center}| 和 \verb|\end{paragraphalign}|。对齐方式可选 \texttt{left}、\texttt{center}、\texttt{right} 或 \texttt{justify}，作用于区块中的普通正文。区块可跨页或正文框，错误指令对应的正文沿用文档默认对齐。

\paragraph{字母与数字间距（0.1.87）}在「正文与段落」应用并保存，可让正文数字随字母按实际笔迹宽度排列。旧项目打开时延用原间距。

\paragraph{独立样本库（0.1.88）}从左下样本库或 Alt+6 打开。采集页的快捷入口将带入当前书写者。筛选结果按字符分组，右侧可切换份次并查看笔迹和来源。展开组后可逐份多选和导出。

\tableofcontents
\mainmatter
\part{认识与开始使用}
\chapter{应用、项目与资源}\label{ch:basics}
\section{启动编辑器}
macOS 用户解压当前交付包，打开完整的 \code{InkStructEditor.app}，无需另装 Python。更新前先保存项目并退出旧版，再打开新应用；不要单独移动包内工作程序。

本册对应 \code{InkStructEditor-macos-arm64-auto-scripts.zip}，包含公式图形化设置重设计，并保留旧版键盘焦点修复。正常打开后，菜单栏属于 InkStruct，左侧源码区可输入。

源码用户可用仓库启动脚本。Windows 需使用配套的启动或离线安装材料，不能运行 macOS 的 \code{.app}。

\section{三类对象要分开保存}
\par\noindent\begin{tabularx}{\linewidth}{@{}P{28mm}P{34mm}Y@{}}
\toprule
对象 & 常见格式 & 用途\\\midrule
文档源码 & \code{.md}、\code{.txt} & 保存要写的内容与标记；单独保存不带字形、评价和完整样式。\\
原生项目 & \code{.inkdoc} & 保存内容、排版设置、已选字形、内嵌资源、模板绑定与评价；日常继续编辑和跨机传递优先使用。\\
可复用模板 & 模板 \code{.json} & 规定纸张、正文区、固定填写区和规则；不包含某份报告的正文。\\
\bottomrule
\end{tabularx}

SVG/PNG 是结果文件，不代替项目。生成目录提供汉字等轨迹，符号 ZIP 提供录制样本；项目内嵌已纳入的资源，不自动复制整个外部字库。

\section{旧项目如何处理}
旧数字表格项目 \code{.inktable} 可导入并另存为 \code{.inkdoc}，引用的样本包须仍在原位置且内容一致。迁移后基线可能变化；严格复现旧排版时保留旧项目与旧工具。

\code{.inkreview} 用于独立字形检查，保存评价与生成目录引用，不内嵌全部轨迹。本册的保存与恢复以 \code{.inkdoc} 为准。

\section{屏幕尺寸与纸面尺寸}
预览缩放只改变屏幕显示，不改变纸面毫米尺寸；“100\%”不是物理校准。普通页面默认为 A4、四边页边距 20 mm；其他纸型和正文区域通过页面模板设置。

\chapter{十分钟完成第一份文档}\label{ch:quick}
本例只使用内置数字、小数点和几何分数线，无需先导入个人字库。随手册提供 \code{examples/01-first-document.md}，可以导入，也可以直接复制下方源码。

\section{输入与观察}
\begin{enumerate}
\item 打开应用，确认\ui{设置 → 正文}中的内容模式为\ui{Markdown 混排}，Markdown 格式为\ui{普通 Markdown}。
\item 点击左侧源码区，输入或粘贴下列内容。独立公式的起止双美元符号分别占一行。
\item 等待底部显示“预览已更新”。右侧应出现数字正文、两列表格和一个分数；允许出现数学校准提醒。
\end{enumerate}
\par\noindent\begin{minipage}{\linewidth}
\begin{Source}
12.3

| 1 | 2 |
| --- | --- |
| 3.4 | 5.6 |

$$
\frac{1}{2}
$$
\end{Source}
\end{minipage}\par\medskip
\shot{workspace.png}{1}{数字、表格与分数的实际编辑界面。左侧为源码，右侧为手写页面。}

\section{调整一项设置}
将光标放在表格源码中，打开\ui{设置 → 元素 → 表格设置…}。先保留总宽 160 mm、最小行高 8 mm、数字高度 3.8 mm、留白 0.3 mm。确认当前作用范围后修改对齐方式，观察预览，再取消或确认。学习时每次只改一个参数，更容易理解效果。

\section{保存并导出}
\begin{enumerate}
\item 点击顶部\ui{保存}，选择项目名称，保存为 \code{.inkdoc}。下次用\ui{打开 → 打开项目…}继续。
\item 点击\ui{查看问题}。错误需要修复；“数学校准尚未人工复核”提醒表示需要在正式书写前检查纸面效果。
\item 点击顶部\ui{导出 SVG…}，选择一个尚不存在的输出目录。多页文档会按页输出。
\item 核对导出目录中的 \code{page-001.svg} 和说明清单。它们是机器输出；项目文件仍用于后续修改。
\end{enumerate}

\note{不要将带中文标题的示例作为“无需字库”的首次测试。中文能出现在输入框中，但预览仍需要相应手写轨迹。首次练习成功后，再按第\ref{ch:resources}章导入个人资源。}

\part{编辑排版}
\chapter{界面与日常编辑}\label{ch:workspace}
\section{常用入口在哪里}
\par\noindent\begin{tabularx}{\linewidth}{@{}P{38mm}Y@{}}
\toprule
入口 & 用途\\\midrule
\ui{文件}菜单 & 打开项目、导入 Markdown/文本、保存源码、打开恢复副本。\\
\ui{文件 → 保存项目} & 保存原生项目；第一次保存时选择位置。\\
\ui{编辑 → 撤销 / 重做} & 撤销或恢复正文编辑。评价有另外的撤销按钮。\\
顶部\ui{设置} & 打开右侧文档设置，分为正文、元素、字形、问题；保留上次选择的分类。\\
顶部\ui{字形评价} & 展开或收起底部评价区。\\
顶部\ui{导出 SVG…} & 输出完整文档的写字机文件。旁边下拉菜单提供参考页面与评价 CSV 导出。\\
底部\ui{查看问题} & 查看当前排版诊断，定位错误或模板问题。\\
\bottomrule
\end{tabularx}

左侧源码、右侧预览是主要工作区。拖动中间分隔线可以分配宽度；不需要的设置或评价面板可以收起。关闭设置面板不会取消已经确认的文档设置。

\section{编辑、定位与撤销}
直接输入或粘贴源码后，预览稍候更新。输入法的候选和未提交拼音不应成为正文。重排时可能暂时保留旧页，等到“预览已更新”后再判断最终结果。

点击预览中的字可定位相应源码并选中具体字形。点击任务方框会改变源码中的勾选状态，支持撤销。图片点击定位整段图片源码，图片本身不作为字形评分。

\par\noindent\begin{tabularx}{\linewidth}{@{}P{31mm}P{40mm}Y@{}}
\toprule
动作 & macOS & Windows\\\midrule
打开项目 & Command+O & Ctrl+O\\
保存项目 & Command+S & Ctrl+S\\
剪切 / 复制 / 粘贴 & Command+X/C/V & Ctrl+X/C/V\\
全选 & Command+A & Ctrl+A\\
撤销 & Command+Z & Ctrl+Z\\
重做 & Command+Shift+Z & 以\ui{编辑}菜单为准\\
\bottomrule
\end{tabularx}

\ui{编辑}菜单的操作发给当前有焦点的输入框。若正在编辑设置中的数字或文字，撤销与粘贴也可能作用于该控件；需要操作正文时，先点击源码区。源码标题旁的\ui{粘贴}按钮会明确把剪贴板文本插入正文光标处。

\section{正文设置}
\entry{设置 → 正文}
\par\noindent\begin{tabularx}{\linewidth}{@{}P{33mm}P{23mm}Y@{}}
\toprule
参数 & 新文档默认 & 调整的效果\\\midrule
字高 & 6 mm & 控制普通文字的标称字高；实际墨迹高度由字形决定。\\
行高 & 9 mm & 控制普通行的基线间距；复杂内容可能需要更高的行。\\
笔宽 & 0.3 mm & 控制统一描线宽度，不模拟原笔压变化。\\
段距 / 标题段距 & 2 / 3 mm & 调节块与块之间的空隙。\\
首行缩进 & 0 mm & 控制普通段落的首行起点。\\
列表 / 引用缩进 & 8 / 4 mm & 调节相应内容的缩进。\\
\bottomrule
\end{tabularx}

通过\ui{标题字号…}设置各级标题的倍率。字号变大可能改变断行和分页，重新检查表格、公式和模板区域。不要用缩放预览代替调整实际字高。

\section{选择内容模式}
\ui{Markdown 混排}解释标题、表格和公式等语法；\ui{纯文本检查}用于按字面文字检查字形。Markdown 格式选择只在混排下有意义。已有项目打开时会恢复自己的模式与设置，不一定与新建默认值相同。

\chapter{编写 Markdown 混排内容}\label{ch:markdown}
\section{正文、标题与列表}
空行分隔段落。标题使用一至六个井号，井号后留一个空格。列表、引用和链接可按下例编写；示例含中文和标点，需先导入相应轨迹。
\par\noindent\begin{minipage}{\linewidth}
\begin{Source}
# 测量记录

这是第一段。下一段前留一个空行。

## 操作步骤

1. 检查接线
2. 记录读数

- 第一次测量
- 第二次测量

> 注意单位一致。

[查看说明](https://example.com)
\end{Source}
\end{minipage}\par\medskip

链接只排版可见文字，不把目标网页带入页面。强调、加粗和行内代码保留语义，但手写资产不会因此自动变成系统字体的粗体或斜体。任意 HTML 与 CSS 不是本编辑器的排版接口；需要控制页面位置时使用模板。

\section{任务列表}
\par\noindent\begin{minipage}{\linewidth}
\begin{Source}
- [ ] 12.50
- [x] 13.25
\end{Source}
\end{minipage}\par\medskip
方框和勾线由可书写折线绘制。点击预览中的方框会更新源码，可使用正文撤销恢复。列表文字仍需要字形资源；方框不需要另外采集。只有当前预览完成后才可操作，过期页面不会修改新源码。

\Needspace{8\baselineskip}
\section{脚注}
\entry{设置 → 元素 → 脚注设置…}
\par\noindent\begin{minipage}{\linewidth}
\begin{Source}
12.50[^1]

[^1]: 0.05
\end{Source}
\end{minipage}\par\medskip
脚注按首次引用顺序编号，放在页底；同一标签多次引用共享注释。长注释可续页。设置可调整标记与注释字号、行高、与正文的距离、分隔线以及页底预留比例。

定义放在顶层。脚注可以包含已支持的代码、公式、表格和图形；复杂结构仍遵循各自的分页规则。不要在定义内递归引用另一个脚注。缺失定义或不可容纳的内容会显示诊断，先修正再导出。

\section{代码与字面内容}
\entry{设置 → 元素 → 代码设置…}
使用单反引号包住行内代码，或使用三反引号围成代码块。代码区不解释数学定界符或普通 Markdown 标记。例如：
\par\noindent\begin{minipage}{\linewidth}
\begin{Source}
```text
12.50
  13.25
```
\end{Source}
\end{minipage}\par\medskip
代码块保留空行、空格和制表位。默认字高 4 mm、字格宽 5 mm、最小行高 6 mm、制表位 4 列。默认超宽行报错；可在设置中明确允许换行，续行会带标记。程序代码中的英文字母和符号也必须有真实字形资源，不自动使用等宽印刷字体作为机器轨迹。

\section{图片与单线图}
\entry{设置 → 元素 → 图片设置…}
先在源码中独占一段写图片标记，再在图片设置中选择这一项并导入本地文件。
图片设置的应用范围可选择文档默认或任意一张图片。默认设置用于沿用默认的图片及后续新增图片，已有独立设置不受影响。选择具体图片后可设置个性化参数，或勾选恢复沿用文档默认；设置及资源随项目保存，取消不写入项目。没有图片时也可以预先配置默认尺寸。
\par\noindent\begin{minipage}{\linewidth}
\begin{Source}
![图一 测量示意](assets/chart.png)
\end{Source}
\end{minipage}\par\medskip
方括号文字作为图题；无文字则不显示图题。图题不自动编号。0.1.24 起默认高 60 mm，宽度按原图比例计算，图题间距 2 mm、居中；默认页不再提供宽度输入。单张图片可在右侧独立选择按高度或按宽度等比设置；旧文档尺寸保持原值。图片与图题整体分页，不支持正文环绕或表格单元格内图片。

\par\noindent\begin{tabularx}{\linewidth}{@{}P{31mm}Y Y@{}}
\toprule
导入类型 & 使用方式 & 可用输出\\\midrule
PNG / JPEG & 作为完整图片，与手写图题一起预览。 & 参考 SVG/PNG；机器 SVG 导出前提示跳过图片，保留占位与图题。\\
受控单线 SVG & 导入作者确认的中心线路径，设置纸面尺寸和曲线容差。 & 校验通过后可与正文一起导出机器路径。\\
模板底图 & 在模板设计器中导入，帮助定位预印纸。 & 仅供参考，不作为要书写的正文图像。\\
\bottomrule
\end{tabularx}

单线 SVG 需要无填充、明确黑色描边和有效宽高；不接受文字元素、外部链接、脚本、滤镜、虚线和任意字体轮廓。曲线按最终纸面误差转换为折线，默认容差 0.05 mm。调小容差可增加细节，也会增加路径点。导入失败时应按诊断修改原图，不把更改扩展名当作格式转换。

\section{Obsidian 静态扩展}
\entry{设置 → 正文 → Markdown 格式 → Obsidian 静态扩展}
\par\noindent\begin{minipage}{\linewidth}
\begin{Source}
[[资料/测量|12.50]]

> [!note]- 测量记录
> 13.25

![[附件/原记录|原记录说明]]
\end{Source}
\end{minipage}\par\medskip
双链只显示标签；无别名时显示目标文字。Callout 按引用样式展开，折叠标记不会把纸面内容隐藏。带显式标签的嵌入只显示标签，并提醒目标内容未展开；没有标签的嵌入会报错。该模式不读取 Obsidian 库，也不执行插件、CSS、动态查询或递归嵌入。

管道表格中的别名分隔符应转义，例如 \code{[[资料/测量\|12.50]]}。需要真正放入图片时，使用普通 Markdown 图片和显式资源导入流程。

\chapter{表格排版}\label{ch:tables}
\section{默认设置与个性化设置}
\entry{设置 → 元素 → 表格设置… → 应用范围}
应用范围可选择“文档默认”或正文中的任一表格，无需先将光标放在表格内。文档默认影响沿用默认的表格及后续新增表格；个性化设置只作用于所选表格。确定后应用于所选范围，切换范围会重新载入该范围已保存的设置。

当前表格可以恢复沿用默认，取消独立设置。已有项目保存自己的尺寸，打开旧项目不会强制套用新默认值。

\section{两种表格源码}
简单表格用管道语法。第二行分隔表头和数据；冒号表示对齐。
\par\noindent\begin{minipage}{\linewidth}
\begin{Source}
| 1 | 2 | 3 |
| :--- | :---: | ---: |
| 12.50 | 13.25 | 14.00 |
| 15.00 | 16.75 | 17.50 |
\end{Source}
\end{minipage}\par\medskip
需要合并单元格时，使用显式闭合标签的 HTML 表格：
\par\noindent\begin{minipage}{\linewidth}
\begin{Source}
<table>
  <tr><th colspan="2">1</th><th>2</th></tr>
  <tr><td rowspan="2">3</td><td>4.5</td><td>6.7</td></tr>
  <tr><td>8.9</td><td>10.1</td></tr>
</table>
\end{Source}
\end{minipage}\par\medskip
\code{colspan} 表示跨列，\code{rowspan} 表示跨行。可使用受支持的单元格水平/垂直对齐和行内公式，不支持嵌套表格或任意 CSS。框线由排版器绘制，相邻公共边不会重复描线。

\section{默认值与列宽选择}
\par\noindent\begin{tabularx}{\linewidth}{@{}P{37mm}P{25mm}Y@{}}
\toprule
参数 & 默认值 & 使用建议\\\midrule
表格总宽 & 160 mm & 先核对可用正文宽度；窄模板需手动调小。\\
最小行高 & 8 mm & 内容可以撑高；启用固定行高后放不下就报错。\\
数字高度 & 3.8 mm & 独立于正文字高和公式字号。\\
四边留白 & 各 0.3 mm & 控制笔迹与框线净空，不建议为压缩表格降到重叠。\\
\bottomrule
\end{tabularx}

列宽可按内容分配、等宽、比例或固定毫米值。固定列宽的合计必须等于总宽。逐列对齐可填“沿用、左、中、右、小数点”，使用逗号分隔。小数点对齐适合小数位数不同的读数；它不会把 \code{001.20} 改写为 \code{1.2}。

当表格放不下时，先查是哪一列内容太长，再调整列宽、数字高度、公式字号或内容组织。程序不会静默压缩字形来填满单元格，也不会自动缩小整表。

\section{跨页与合并单元格}
普通行按可用页高分页，表头可重复。跨行合并连接的若干行必须作为一个完整行组同页排版，不会拆开单元格内部内容。如果整个行组高于完整正文区，继续增加页数也无济于事；需要拆分合并、拆表或调整正文区与字号。

重复表头属于原内容的重复展示，不是新输入的数据。检查分页时重点观察最后一行、续页首行以及长公式所在行，避免仅看第一页就导出。


\chapter{公式输入与手写样式}\label{ch:math}
\section{输入规则}
0.1.24 起可在公式中直接使用汉字，包括下标与上标，例如 \code{U_{模块}}、\code{U_{万}}、\code{温度}；分式、根式和矩阵同样支持。汉字使用已载入的真实字体，缺字继续定位到原文，不自动改用屏幕字体。\code{\text{模块}} 写法仍然支持。
Markdown 行内公式使用单美元符号；独立公式的起止双美元符号分别独占一行。代码区内的美元符号保持字面内容。
\par\noindent\begin{minipage}{\linewidth}
\begin{Source}
12.50 的一半为 $\frac{12.50}{2}$。

$$
\sqrt{12.50}
$$
\end{Source}
\end{minipage}\par\medskip
示例中的中文需要汉字资源。公式命令名不会逐字书写，但其显示结果需要实际字形，例如 \code{\alpha} 需要相应希腊字母轨迹，\code{\sin} 需要函数显示所用字母。内置数字不能代替完整数学字库。

\section{常用结构速查}
\par\noindent\begin{tabularx}{\linewidth}{@{}P{33mm}Y@{}}
\toprule
结构 & 输入示例\\\midrule
上下标 & \code{x_1^2}、\code{x^{12}}\\
分式 & \code{\frac{1}{2}}、\code{\dfrac{1}{2}}、\code{\tfrac{1}{2}}\\
根式 & \code{\sqrt{12}}、\code{\sqrt[3]{8}}\\
伸缩括号 & \code{\left(\frac{1}{2}\right)}\\
上下限 & \code{\sum_{i=1}^{n}}、\code{\int\limits_0^1}\\
文字与单位 & \code{\text{电压}}、\code{\mathrm{m^2}}\\
修饰符 & \code{\vec{x}}、\code{\overline{x}}、\code{\hat{x}}、\code{\dot{x}}\\
\bottomrule
\end{tabularx}

上下标超过一个字符时加花括号。\code{\left} 与 \code{\right} 要配对；不可见一侧可写点号。管道表格中的公式尽量用 \code{\vert} 或 \code{\Vert} 表示竖线，避免未转义的管道符被当作分列标记。

公式作为整体参与换行与分页，不会自动在中间折行。超宽时，可改为显式多行结构、减小公式字号或扩大有效区域。上下标层级过深、字形过小也会报错，不会无限缩小到不可书写。

\section{在设置窗口调节结构}
\entry{设置 → 元素 → 公式设置…}
设置窗口按结构分类：字号与字形、间距与数学轴、上下标、分式、根式、积分与大运算符、矩阵与向量、括号与定界符、多行对齐、公式编号、修饰与向量。

左侧选类别，在上方紧凑结构图旁的小输入框直接调节参数。局部短箭头标明参数对应的尺寸，聚焦时突出显示；窗口放大不会拉开字符与控件。示意图使用正体字形，保持固定比例，数值变化请以下方预览为准。下方实际字形预览使用当前项目笔迹，有效数值输入后约 0.1 秒刷新，无需按回车或离开输入框。预览倍率固定，修改字号不会被自动缩放抵消；可在倍率下拉框主动缩放，较大公式可滚动查看。超范围输入显示原因，不继续显示旧预览。

示例输入只填公式内容，不加美元符号；可切换行内与独立预览。示例下拉框提供小矩阵、带框线数组、连分式等试排内容，切换分类会保留各类自定义示例。示例不会自动插入正文；缺字仍显示红框。

默认公式字号为 6 mm，与表格数字高度独立。字号可在 mm 与 pt 之间联动显示；百分比参数以当前层公式字号为参照，例如分子、分母和上下标可能已经缩小。将鼠标悬停在参数旁的\ui{i}上，可查看参数含义和按主字号估算的毫米尺寸；点击图标也可打开提示。分类和预览标题旁的\ui{i}提供结构规则与资源说明。估算不能代替实际嵌套尺寸。勾选\ui{显示辅助线}可查看灰色基线和蓝色数学轴；所有尺寸箭头和辅助线均不进入机器文件。

参数框旁的回转箭头只恢复单项；\ui{当前分类恢复默认}重置该类尺寸、来源及该类对齐选项。\ui{应用设置}应用到文档并随项目保存，\ui{取消}保留原设置。本版没有逐条公式的独立样式覆盖，调整时应同时查看文档中其他公式。

普通上下标不再提供“净间距”手动项。程序按实际笔迹边界计入笔宽，自动保持水平留白及上下标之间的必要间隔；上移、下移设置仅控制最小位移，必要时会自动避让。分类标题旁的\ui{i}说明这一行为。

在\ui{积分与大运算符}中，\ui{上下限字号}独立于普通上下标字号，作用于求和、连乘、积分等的右侧或堆叠上下限。比例以当前层公式字号为基准，默认 70\%，可调范围为 50\%--100\%；上下限内部的嵌套上下标仍沿用普通上下标比例。

旧项目缺少独立上下限比例时，继承原有上下标比例后独立保存。普通上下标按新的自动留白规则重新排版，旧净间距值不再控制几何；打开旧文档后请检查分页与公式位置。

\section{使用录制的减号、根号与括号}
三类构件可以分别选择\ui{几何构件}或\ui{录制字形（缩放适配）}。自 0.1.27 起，新文档分数线默认使用字体中的数学减号横向拉伸，保留纵向笔迹起伏；根号和括号仍默认几何构件。旧文档保持已保存来源。缺少减号时可补录，或明确选择几何构件，程序不会静默替换。
\begin{enumerate}
\item 在\ui{分式}中选择录制来源，用数学减号的轨迹按内容宽度延长分数线。
\item 在\ui{根式}中选择录制根号，程序适配钩部、竖部与上横线；根内内容变长时主要延长横线。
\item 在\ui{括号与定界符}中选择录制来源，对相应括号按内容高度与构件宽度缩放。矩阵、分段和编号括号也受这一选择影响。
\end{enumerate}
\shot{settings-fraction.png}{0.94}{紧凑分式设置：短箭头连接附近输入框，说明收纳于 i 图标，下方核对实际笔迹。}

数学减号是 U+2212，不等同于正文中的 ASCII 连字符。缺样、退化轨迹或不能可靠识别的根号会报错。可补充合适样本，或明确切回几何模式；程序不会悄悄换来源。适配只作用于派生结构，不改原始采集数据。

\section{矩阵、分段与多行推导}
普通矩阵支持 \code{matrix}、\code{pmatrix}、\code{bmatrix} 等环境；\code{cases} 用于两列分段。矩阵以 \code{&} 分列，以双反斜线换行；普通矩阵各行列数应一致，最多 8 行、8 列。
\par\noindent\begin{minipage}{\linewidth}
\begin{Source}
$$
\begin{pmatrix}
1&2\\
3&4
\end{pmatrix}
$$
\end{Source}
\end{minipage}\par\medskip
多行推导可用 \code{aligned}，对齐点放在等号之前。下例需要加号、等号和数学减号等相应资源。
\par\noindent\begin{minipage}{\linewidth}
\begin{Source}
$$
\begin{aligned}
12+3&=15\\
15-2&=13
\end{aligned}
$$
\end{Source}
\end{minipage}\par\medskip
\code{gathered} 逐行居中；\code{split} 用于一个展示公式中的共同对齐点。多行结构仍整体分页，不跨页拆行。需要更多可独立分页的步骤时，应拆成多个独立公式。

\section{编号与引用}
普通独立公式默认不编号。\code{equation} 为整块生成编号，\code{align}、\code{gather} 为每行生成编号；对应星号环境不自动编号。
\par\noindent\begin{minipage}{\linewidth}
\begin{Source}
$$
\begin{equation}
\frac{1}{2}\label{eq:half}
\end{equation}
$$

$\eqref{eq:half}$
\end{Source}
\end{minipage}\par\medskip
\code{\tag{12.3}} 指定可见标签，\code{\tag*{12.3}} 不带括号；\code{\notag} 取消该行编号。\code{\ref} 显示标签文字，\code{\eqref} 显示带括号的引用。引用命令要放在数学定界符内。

\code{\label} 的键全篇唯一，引用可以出现在被引公式之前。重名、缺失目标或在无编号行登记标签会报错。插入前面的公式后，自动编号与引用会更新；不需要手动修改正文数字。

\section{与完整 LaTeX 的区别}
编辑器接受受控数学语法，不编译完整 \code{.tex} 文档。不要粘贴 \code{\documentclass}、\code{\usepackage}、宏定义或文件输入命令。\code{array} 支持 \code{l/c/r} 列及受控框线，但不支持任意列格式或 \code{\multicolumn}。未知命令会定位报错，保留原文供修改。

数学校准提醒表示当前参考线还需实际样张复核；它不同于“缺少字形”的阻断错误。资源齐全后，也应检查上下标、分数线和括号在实际笔宽下是否清楚。

\chapter{设计和使用通用页面模板}\label{ch:templates}
\section{模板库、本文布局与默认设置（0.1.35）}
简单模板以一张纸的宽高和四向边距为布局，正文自动换行分页。字号、行距和段距属于排版样式，从齿轮的“正文与段落”设置。自定义模板通过起始页组和循环页组组织源布局；循环的是布局，实际页码连续递增。

主窗口\ui{模板库}用于管理可复用模板，显示缩略图、页面版式、页组顺序和字段声明。导入保存独立库内副本。\ui{保存修改}更新当前库内模板并选中最新修订，旧修订可在\ui{显示历史修订}中查看。编辑待入库的最近文件时，保存修改会写回原 JSON。\ui{另存为新模板}创建独立副本。移出条目可恢复，已保存的文档使用其内嵌模板。

排版右侧\ui{文档资源}用于更换本文字体、编辑本文布局、更换模板、填写字段和定位实际适配问题。更换模板先试排确认，保留正文与冻结资源；编辑菜单统一撤销重做。新文档默认仍在左下角设置配置，齿轮里的模板库只提供跳转入口。

\section{工作台与字段填写}
自定义工作台左侧维护起始、循环与未使用布局，中间画布显示正文编号及连接，支持选框、拖动移动和右下角缩放，右侧显示所选对象属性。共用布局的引用与独立副本明确区分。简单模板只呈现纸张、边距和预览。草稿自动收集有效输入，没有各页的“应用草稿”按钮；无效输入保留并提示，取消整个工作台不修改外部文档。确认按钮明确为“应用到当前文档”或“保存为新模板”。

字段表单只显示模板声明的字段，确定后写回单份 Markdown，保留未改内容和换行。重复定义及语法错误需回主编辑器定位处理。布局结构检查与真实内容试排分开，试排使用独立资源解析器，不改变当前选样。旧模板继续保留原设计器与源码块绑定，迁移为显式副本操作。

\section{有序正文框、固定字段与 PDF}
每页可有多个正文框，按明确编号连接。页面顺序和本页框编号共同决定唯一正文流，画布显示编号与连接关系，不猜空间阅读顺序。字段框引用稳定标识，多处引用重复显示，不能触发续页；页码使用独立自动框。固定框放不下时报告溢出，不默认缩字或丢弃内容。

PDF 支持按指定页码顺序一次导入多页背景。默认采用首个所选页的物理尺寸，可确认或保留原纸张；修改纸张时背景分别按宽高拉伸铺满，文本框保持毫米值。PDF 和图片均只作为参考，不自动描摹。预印纸只写内容；空白纸只输出明确的框线和内容。

\section{单份 Markdown 命名字段}
整份文稿只有一份源码。字段标识与显示名、具体框名分开；多个框可引用同一标识。字段块外全部进入连续正文，字段文字不在正文重复。示例：
\begin{verbatim}
:::field name
张三
:::

:::field student_id
20260001
:::

# 实验目的
这里开始是连续正文。
\end{verbatim}

重复定义、必填缺失或为空、未声明字段、缺少结束标记和嵌套字段块都阻止正式导出，并可定位源码。缺少结束标记不吞掉后续正文。代码围栏中的标记按普通代码处理。可选字段缺失时留空；字段支持已支持的 Markdown，但整体不得跨出固定框。

\section{从印刷纸样识别正文框}
导入图片或 PDF 后，可直接在底图上调整自由、横线与方格正文框，增删区域并核对行列。方格文字逐格放置，换行后按框和循环页顺序继续。保存模板后再次使用时沿用已确认的区域。

\section{结束、补齐与适配}
末尾补齐循环页组默认关闭，正文在哪一页结束就在哪一页停止。开启后只补完已经进入的当前循环组，补齐页仍绘制固定信息。空正文或仅字段时输出起始组一次；没有起始组则输出首循环布局，勾选补齐时输出完整首循环。已填写但没有实际框承接的内容明确报错。

正文跨框传递未消费内容与排版状态，按新宽度断行，不截断 Markdown 重解析。表格沿用原子行组和重复表头，公式和图片遵守原有尺寸规则。完整循环无法取得进展时停止并定位，页数保护与导出检查保留。

模板 v3 与原生项目 v3 显式版本化。旧版不自动迁移；复制旧几何时记录源版本和哈希，保留旧文件，原绑定、锚点和倍率需人工核对，不根据名字改写正文。以下各节记录旧模板的兼容操作。

\section{兼容旧模板：内容与绑定}
模板规定“在哪里写”，Markdown 规定“写什么”，绑定规定“哪一块内容进入哪个填写区”。同一模板可以用于不同文档；同一文档也可换用模板并重新绑定。学校报告、记录单或自定义表单只是不同配置，不是应用固定的栏目。

\par\noindent\begin{tabularx}{\linewidth}{@{}P{32mm}Y@{}}
\toprule
对象 & 用法\\\midrule
正文区 & 未绑定到固定区的正文按顺序流入，可自动续页。\\
固定填写区 & 姓名、记录编号、结论等完整内容块的指定区域，不能内部跨页。\\
锚点 & 定义首字横坐标与基线，帮助文字对齐预印下划线。\\
规则线 & 空白纸上需要绘制的框线或折线；与正文分开安排执行顺序。\\
底图 & PDF/图片转换来的纸面参考，不自动描摹为书写路径。\\
内容字段 & 为填写区附上稳定编号、显示名称、类型和必填规则，便于复用及 Agent 操作。\\
\bottomrule
\end{tabularx}

\section{新建一个中性模板}
\entry{设置 → 元素 → 页面模板… → 新建 / 设计模板…}
\begin{enumerate}
\item 确定真实纸张宽高。默认起点为普通 A4，不自动套用某所学校的表单。
\item 切换首页和续页，设置各自正文区。坐标原点在纸张左上角，横向向右、纵向向下，单位毫米。
\item 用\ui{绘制填写区}建立固定区；用\ui{放置锚点}设置首字基线；用\ui{绘制直线}添加需要书写的规则线。
\item 在属性表核对名称、位置和尺寸，修改数值后点击\ui{写入}。未写入的修改需先处理，再切换对象或确认。
\item 检查几何预览与问题列表，通过后导出模板 JSON，回到页面模板窗口绑定内容。
\end{enumerate}

当前每种页型只有一个正文流区域；首页和续页正文等宽，续页正文高度不少于首页。正文与固定区不能重叠。暂不支持多栏、不同栏宽的自动重排或任意页型序列。

\section{画布操作与精确调整}
\ui{选择 / 移动}工具可拖动区域、锚点或整条折线，区域四角可改尺寸，折线顶点可单独移动。每次拖动松开形成一次撤销操作，Esc 可取消本次拖动。多选后可使用对齐与等距分布；排列前核对选中的对象。

默认吸附间距为 1 mm，可修改或关闭，按住 Alt 可临时关闭吸附。标尺、缩放与底图透明度帮助观察，但不改变毫米坐标。拖动导致越界、重叠或违反页型约束时不会接受该结果；数值/JSON 草稿的错误通过问题列表修正。

给区域命名后再绑定内容。已在使用的名称需要修改时，优先使用外层页面模板窗口的\ui{联动重命名填写区 / 锚点…}，它会一起更新当前项目中的引用。直接改 JSON 名称不会自动推断原来的绑定关系。

\section{使用 PDF 或图片作参考}
设计器中的\ui{导入当前页底图…}支持本地 PNG/JPEG/PDF。PDF 导入时先选择页码，再确认目标是首页还是续页；导入 PDF 第 2 页与选择“续页”是两件事。

\begin{description}
\item[仅替换底图] 保留当前纸张和区域。若 PDF 尺寸不同，参考图会按当前纸张铺满，可能改变比例；需认真核对。
\item[按 PDF 本页尺寸新建布局] 按页面物理尺寸建立通用正文区，替换旧几何、底图和偏移。原有绑定需重新检查。
\end{description}
PDF 不会自动识别姓名栏、框线或正文。应使用底图手动设计，或让 Agent 提供经过验证的模板。导入可取消；超时可改用更简单的 PDF 或预先准备的图片。

\section{绑定固定字段：完整示例}
设想一张通用记录单：顶部留“记录编号”和“日期”，下面是自动续排的正文。先在首页画好两个固定区，将正文区放在它们下方。然后在源码中让编号、日期和正文各自成为独立段落。
\par\noindent\begin{minipage}{\linewidth}
\begin{Source}
20261014

2026.10.14

12.50

13.25
\end{Source}
\end{minipage}\par\medskip
\begin{enumerate}
\item 回到\ui{页面模板…}，选择记录编号填写区和第一段源码块。
\item 按需选择首页/所有页、倍率和基线锚点，点击\ui{绑定所选源码块}。
\item 将第二段绑定到日期区；后面未绑定的两段继续进入正文流。
\item 确认页面模板设置，等待实际字形排版，检查固定区是否越界。
\end{enumerate}
同一源码块不能同时绑定到两个固定区。解除绑定后，该内容回到正文；不会删除源码。删除或整体替换一个绑定块后，绑定可能失效，应重新选择块或解除旧绑定。

若要每页重复，首页和续页都需有同名填写区及所用锚点。固定区内容不能靠自动截断或压扁字形塞入；倍率也不会同比缩小表格的独立毫米尺寸。

\section{为复用模板声明内容字段}
在设计器的\ui{内容字段…}中新建字段，选择填写区，填写稳定编号和用户可读名称；设置内容类型为文字、公式、表格或任意完整块，并决定是否必填、在哪些页面出现。

使用已声明字段时，页面策略、倍率与锚点由模板控制，需要回到设计器修改。每份文档仍显式选择自己的源码块，程序不会因为字段名叫“姓名”就自动寻找文字。

添加语义字段会使用模板 v2；旧版应用可能无法读取，保留原模板副本。模板声明不保存某份报告的人名或正文。删除字段声明也不会自动删除原绑定，清理时应分别检查。

\section{预印纸、空白纸与试写}
\par\noindent\begin{tabularx}{\linewidth}{@{}P{26mm}Y Y@{}}
\toprule
模式 & 页面预览 & 机器输出\\\midrule
预印纸 & 显示纸面底图，或灰色参考框线。 & 只写内容，不重复画印刷底图与参考框线。\\
空白纸 & 显示配置的规则线与内容。 & 先画显式规则线，再写各固定区与正文；不复制扫描图上的印刷标题。\\
\bottomrule
\end{tabularx}

先用少量定位内容试写，测量整体横纵偏移，再在页面模板中设置毫米校准偏移。预印模式移动书写内容、保持纸面参考不动；空白模式规则线与内容一起移动。当前仅支持平移，不校正旋转或纸张缩放。

\section{查看适配问题}
\entry{查看问题 → 适配报告… / 修正模板…}
等待最新预览完成后打开适配报告，检查必填字段、绑定状态、缺字和溢出。选中问题可定位具体填写区或源码；也可另存 JSON 供 Agent 处理。

“模板几何有效”只说明区域合法，不代表正文能放下；“内容可导出”也不代表已完成设备校准。改动正文、字号、模板或资源后，要重新查看最新报告。使用\ui{普通页面}会移除模板并清除固定区绑定，内容返回普通正文流。


\section{连续正文流中的分页与换框}
多文本框连续正文流中，\code{\framebreak} 跳到下一个编号正文框；同页下一框不增加页码。\code{\newpage}、\code{\pagebreak}、\code{\clearpage} 跳到下一物理页，跳过当前页剩余正文框。普通 Markdown 页面也支持分页。

命令可独占一行或位于段内；连续换框逐个推进，空白页开头的分页不额外造页，尾随控制不额外创建空白输出页。列表标记随正文进入新框。固定字段、表格单元格和图片说明不执行分页或换框；普通页面及旧版固定区不执行换框。\code{\pagebreak[4]} 强制分页；强度 0--3 不强制推进。控制命令写在公式外，原子公式本身不拆分。内容由 AI 根据 Markdown 规范生成后复制或导入，不增加插入控制符的按钮。旧语法项目打开时显式迁移换框命令，保留原文件并另存。

\chapter{字形资源与质量评价}\label{ch:resources}
\section{本机大容量字体（0.1.118）}
\entry{主窗口 → 字体库 → 制作字体}
预计候选超过 30,000 份时，制作前会提示。选择\ui{保留全部，继续制作}可保留原定类型和全部候选。候选较少但含大量来源详情时，也可主动勾选\ui{使用本机大容量模式}。导入本机大容量包同样需要确认。大型字体的制作、校验和导入可能占用较多内存、磁盘空间与时间。
字体库会标明\ui{本机大容量}。需要跨软件使用时，从\ui{导出字体包}选择\ui{通用 v1 副本（选择候选）}。可以手动勾选，或先为每字每种字形各选一份，再逐项调整。导出的副本依通用规格校验。

\section{生成后的评分}
在生成页勾选来源分、多样性分或两者。「计算方式」也控制评分进程。自动模式在大批次使用最多 8 个进程，实际数量受处理器与可用内存限制。评分阶段单独显示已评分字数，取消后仍可使用已生成的字形。

\section{导入统一字库（0.1.117）}
\entry{主窗口 → 字体库 → 导入字体}
选择 InkStruct 或 OpenInk 的通用 v1 字库后即可排版。导入旧 InkStruct 字体包时，程序在后台校验并转换，保留原包供旧文档找回原来使用的字形。默认转换包含字形来源和数学轴。需要保留逐点压力和时间时，使用导入按钮下拉菜单中的\ui{导入旧字体并携带压力与时间}。转换失败时，已通过校验的原包仍可使用。

\section{字体库与当前文档（0.1.35）}
\entry{主窗口 → 字体库}
字体库左侧搜索与筛选，右侧分别查看真实笔迹、字符和多份样本、来源与版本、文件信息。缺字显示红框，不用系统字形替代。相同字体身份的版本集中显示，文件较新不代表推荐版本。单击预览与双击详情均不自动更换本文。

\ui{用于当前文档…}先验证资源并进行实际试排，列出缺字、来源和冻结选样保留数量。确认才改变本文，支持撤销重做；检查期间文档或资源变化时拒绝旧结果。核心记录、度量和来源未变的冻结选样保持，其他样本明确重新选择。确认不会解除导出检查。

\ui{设为新文档默认}只影响以后新建的文档。\ui{更多操作}可修改显示名称及移出或恢复列表，不修改字体包和旧文档。导出可选用户版或完整副本，原包保留，不能覆盖已有文件。

\section{较大字体（0.1.46）}
可导入展开总量不超过4 GiB的字体，保留单成员32 MiB、每候选点数、成员数和全部哈希/来源检查。展开总量按ZIP内未压缩数据计算：磁盘文件不大，也可能超过旧版上限；项目保存与排版仍有独立容量保护。

user\_001同字最多8份并包含相同英文、希腊、数学等字符的实测用户字体，全结构库包约611 MB，简体支持字集包约454 MB，压缩体积减少约25.7\%。简体支持集为2013年\ui{通用规范汉字表}与当前结构库的6866字交集；结构库缺字仍列明，不能当作全部简体字已经支持。这是字体体积示例，不是应用安装包体积。

两包展开后都超过旧2 GiB限制，需新版程序读取；升级不自动导入新包或替换旧项目资源。先在字体库选择\ui{用于当前文档…}检查实际影响，确认后再应用。字体能载入不等于来源或数学轴已经人工审阅，原有警告保留。

\section{兼容外部生成目录}
\entry{排版 → 文件 → 高级资源（仅当前文档）}
\begin{enumerate}
\item 点击\ui{选择生成字形目录…}，选择包含 \code{trajectories/} 的生成结果目录，供汉字等内容使用。
\item 在字体库的\ui{制作字体包}中选择同一书写者的生成批次、扩展采集目录或 ZIP，核对清单后按类别入包。
\item 等待预览更新，检查缺字；同一个视觉符号可能对应不同 Unicode 字符，应按诊断补齐实际需要的资源。
\item 保存为 \code{.inkdoc}，保留文档实际使用的资源与选样。
\end{enumerate}

内置单线数字和小数点用于快速试排。它们能完成数字练习，但不表示某位书写者已经录制了这些字符。更换资源后应检查整体笔迹是否一致，不将内置数字误认为个人样本。

\ui{刷新当前外部资源}用于重新读取外部生成目录的变化，程序不会自动监控目录。\ui{载入全部字}用于字库检查，会把生成批次的字符载入检查正文；正式报告操作前先保存，建议另开检查项目使用。

\section{缺字与选样}
不同字符和不同样本各有身份。同一字符可有多个录制写法，排版按样本池选择，并保留已有实例的选择；修改附近正文、重排或重开项目不应让未改内容任意换字。

界面能显示某个字符并不说明存在它的轨迹。红色占位或缺字诊断表示资源不够；不要把占位当作成品。补资源后重新排版；若改的是资源版本，应再次观察实际字形。

对整份文档检查补采需求，可让 Agent 使用附录\ref{ch:agent}的资源检查命令。符号采集 CSV 默认每个缺失字符 5 次，汉字生成需求和不支持字符另列在报告中。空缺口清单仅说明资源存在，不证明根号结构、数学轴或页面尺寸已经合格。

\section{字符校正与旧评价}
在预览中点击字符，右侧\ui{字符校正}可调整位置和缩放。选择应用范围后校正可随文档保存，也可在支持的范围另存校正字体。旧项目中的评价记录继续随项目读取和保存。

\chapter{导出、保存与恢复}\label{ch:export}
\section{选择正确的出口}
\par\noindent\begin{tabularx}{\linewidth}{@{}P{33mm}P{39mm}Y@{}}
\toprule
目的 & 入口或文件 & 需要知道的区别\\\midrule
继续编辑 / 跨机传递 & 保存项目，\code{.inkdoc} & 包含内容、设置和内嵌资源；优先作为工作母版。\\
只交换文字 & 保存 Markdown / 文本 & 不带完整字形、评价和样式。\\
查看整页外观 & 更多导出选项 → 导出参考页面… & 输出 SVG/PNG；位图参考页不能当作机器文件。\\
发送写字机 & 顶部导出 SVG… & 完整文档逐页导出，仅可书写路径。\\
分析字形质量 & 更多导出选项 → 导出评价 CSV… & 输出评价和来源，不修改生成器。\\
\bottomrule
\end{tabularx}

页面与机器导出选择尚不存在的新目录，避免覆盖已有结果。多页文件应按页序核对；不要只传第一页。写字机输出以 \code{page-001.svg} 等命名，\code{plotter.json} 记录路径来源、优化距离及相关说明。

\section{写字机 SVG 的含义}
机器页面中的绘制元素仅有 \code{path}，路径指令仅有绝对 \code{M} 和 \code{L}。字体、文字元素、图片和变换矩阵不会留给设备解释。\code{M} 定位到笔画起点，\code{L} 沿折线书写；具体抬笔、速度和力度由设备软件决定。

导出前依次确认：最新预览已经完成；没有阻断错误；字符与图形都具有可书写路径；页面与模板模式正确。缺字、越界和不可支持内容不能通过“忽略提醒”变成完整成品。

含普通位图时，点击导出会说明这些图片不会进入 SVG，包括尚未载入的本地 PNG/JPEG。弹窗分别显示两类数量，展开详情可查看文件名。确认后保留当前图片占位与图题，导出其余文字及受支持的单线图路径；取消不产生输出。原文档及资源不变，\code{plotter.json} 记录原因、页码、源码身份和占位；缺失图片哈希为 null。预览和适配报告继续显示资源缺失错误，整份文档仍未完整可导出。缺失 SVG、未知类型、非法资源、缺字、公式错误与越框仍须修正。需要图片本身参与书写时，应准备经过确认的单线图资源。

\section{比例校正与输出画布}
先严格校验未缩放排版，再以左上角为原点同步缩放输出坐标、SVG 物理宽高及 viewBox，笔宽保持毫米值。例如横纵均为 132.293\% 时，输出画布宽高都乘以 1.32293。文档纸张、预览、分页及冻结选样不改变。选定页独立变换，文件沿用原页码；\code{plotter.json} 记录原始及输出尺寸与变换策略。

缩放前的越界等阻断错误不能通过缩小或筛选规避。缩放后仍校验画布边界、路径及单边 2000 mm 上限；笔宽不随缩小改变，因此贴边笔迹缩小后仍可能被拒绝。定位短线坐标随比例变换，笔宽仍为 0.01 mm。

\section{优化空行程}
\entry{导出 SVG… → 优化路径（允许反转笔画执行方向）}
默认开启。开启后，程序调整路径执行顺序，并允许反向执行某条路径，以减少抬笔移动距离；反转不是另一个可取消的选项。排序和反转不改变路径形状及纸面位置，也不改写原始采样文件。

模板规则线、固定字段和正文等执行区间仍按既定顺序处理，不会为了缩短空行程任意混在一起。清单中的距离表示几何估计，不能直接当作实际节省的时间；落笔方向对具体笔具的影响仍以试写为准。

第一次使用某组样本或笔具时，可以分别导出开启/关闭优化的两组文件比较。两组输出用不同目录，并保留当时的项目版本。

\section{保存项目与多人协作}
\ui{保存项目}保留源码、样式、资源、模板、绑定、选样和评价。\ui{另存项目…}用于试验布局或保留不同版本；\ui{保存 Markdown / 文本…}只输出源码。

从 0.1.82 起，Ctrl+S 在后台保存 Markdown 项目。保存期间可以继续编辑。无改动时跳过写盘；若已有新改动，多次保存请求会合并，当前保存完成后再保存最新内容。

同一项目在另一窗口或另一台机器被修改后，保存可能报告冲突。先另存自己的副本，再比较内容，不反复强行覆盖。编辑器的冲突检测不是多人实时协作或合并系统。

正式更新已有项目时会保留 \code{.previous.inkdoc} 副本。跨机传递后，应在目标机打开项目并重新查看预览；不要只看文件存在就认为资源齐全。复制原始 Markdown 时，还需另外提供字形与图像资源。

\section{打开恢复副本}
\entry{打开 → 打开恢复副本…}
独立应用按会话写自动恢复副本；macOS 默认位置为：
\par\noindent\begin{minipage}{\linewidth}
\begin{Source}
~/Library/Application Support/InkStruct/editor_recovery
\end{Source}
\end{minipage}\par\medskip
源码运行方式的恢复目录是启动工作目录下的 \code{artifacts/editor_recovery/}。每个会话独立保存，重开同一项目或多个窗口不会共用一个草稿文件。

恢复后项目处于未保存状态，核对正文与设置，再选择正式位置保存；恢复副本不会自动覆盖原项目。遇到损坏项目，打开失败应保留当前编辑内容，可继续另存。超限或缺字文档可以保存草稿，但仍需修正后才能导出成品。

\section{交付一份可继续处理的文档}
建议交付原生项目、机器输出目录，以及简短的纸张/笔具/模板模式说明。模板需要复用时另附模板 JSON。若只交付 SVG，接收者可书写，但难以重新调整正文、字体样本或模板绑定。

\section{后台 SVG 导出（0.1.38）}
确认排版在输出位置选择前冻结，导出弹窗中的页码、表格内容、路径优化、横纵比例和定位点设置共同作用于同一份冻结快照。后台期间可继续编辑和滚动，不会改变本次输出；每个窗口同时只运行一次导出。

多页路径优化最多使用四个进程，并受处理器和可用内存限制；单页、小任务及缓存命中沿用串行或缓存路径。进程失败会对未完成页串行重算，不放宽几何和来源检查。

进度窗口可取消，关闭进度窗口也请求取消。取消后清理临时文件，不发布半成品；最终发布已完成时保留成功输出。退出应用须先等待导出取消和清理完成，再次关闭。已有输出目录不能覆盖。

\part{字形生成}
\chapter{从采集样本生成字形}\label{ch:generation}
\entry{统一应用顶部\ui{生成}标签页。默认选择\ui{v6.2 exp（推荐）}；另可选择\ui{v6.1 exp}、\ui{v4 adaptive}。}

\section{准备作者与样本}
填写与采集文件相同的作者 ID，选择该作者的汉字单字 JSON 目录。在生成页直接选择样本目录；采集输出目录可在采集页状态栏查看。目录中应有真实采集的 \code{handwriting_single_v1} 文件；生成轨迹、字体包、编辑字形不能代替真实采集样本建档。数字、字母及数学符号通过字体库直接加入排版资源，不进入本页的汉字部件建档流程。

\section{导入、建档、生成}
\begin{enumerate}
\item 点\ui{导入样本}。程序核对作者、轨迹和动态字段，显示接收数量及跳过数量；详细原因保存于导入记录。没有有效样本时，不能继续建档。
\item 选择生成方式，点\ui{建立档案}。v6.1 自动准备完整来源轨迹及标准笔画映射；v4 准备自适应风格档案。无需另建副档案。部分样本因笔序或映射不可靠而不能成为生成来源，建档记录会列出原因。
\item 在\ui{生成内容}输入汉字，或勾选\ui{生成结构库全部汉字}。点\ui{开始生成}，查看进度。可以取消；串行在当前字结束后停止，多进程停止投递并回收本任务工作进程，保留已核验提交的字。
\item 在生成页点击\ui{制作字体…}，在弹出的制作窗口中为结果命名，作者与生成结果已经带入；按需要加入同一作者的符号样本。制作成功后回到字体库并选中新字体，再选择\ui{用于排版}。
\end{enumerate}
未满足条件的操作按钮会禁用。缺少导入索引时会提示先导入；切换方法后，没有对应档案就需要建档。正在操作时作者和方法暂时锁定。采集目录新增文件后，重新导入再建档；应用不会悄悄更新已生成的结果。

\note{每次导入、建档和生成均保存为独立批次。再次点击生成会建立新批次，不覆盖旧样本、档案或结果。本版全库任务支持取消并保留完成部分，尚未提供应用内断点续跑入口。}

\section{生成方法}
\textbf{v6.2 exp} 在 v6.1 基础上增加竖弯钩与止部关键笔语义准入，拒绝不兼容来源后重新检索，保持真实笔迹及原几何检查。未找到来源时仍明确提示参考替代；来源完整不等于人工质量合格。user\_001 验证合并字体仅重生孔、正，孔的末笔保留参考缺口，其余字体候选和旧文档保持原样。
\textbf{v6.1 exp} 使用经过验证的参考修正、完整部件候选筛选、禁止拆分与位置变体规则、受限关键笔替换及作者单笔回退。保留作者来源点列、供体、变换、拒绝原因和缺口。关键笔替换仅限当前已验证的羊部件第六笔规则，不是任意笔画的自动修复。程序从本次档案重新寻找供体，不会载入另一作者的历史生成字冒充新结果。

本版沿用跨字生成规则：目标字本身的样本不作供体，暂不启用同字原稿优先复用。历史人工冻结的供体和布局只适用于相应档案，不会自动套用到新作者。来源覆盖率不能直接解释为允许同字复用时的生产覆盖率，也不等于人工评价合格率。

\textbf{v4 adaptive} 使用固定的自适应部件融合、强度 1.0 和原有生成路径。它没有完整可靠的逐笔作者来源证明，字体中按来源未知提示。应用不提供神经生成、训练或 CUDA 检测入口。两种方法均不承诺重复运行能产生经过验证的自然随机变体。

详细规则可悬停生成方式旁的\ui{i}提示查看，无需调整部件强度等内部参数。

\section{来源缺口、失败与复核}
作者来源不足但参考几何有效的字可以生成和排版，明确标记参考替代笔序及原因。字体应用时按文档实际使用的字汇总警告；允许继续使用，不强制立即补采。任务完成、作者来源完整、协议有效和人工评价是不同状态。非有限坐标、损坏轨迹或来源校验失败的字不会被当作成功字发布。

生成页显示简短错误和诊断记录位置，不直接显示异常堆栈。部分字失败时，结果目录保留失败记录，字体只纳入完成的字；取消后未生成部分不被当成完成。发布字体前应核对完成数与请求数。

保存为 \code{.inkdoc} 会冻结实际候选、轨迹、度量和来源状态。重开不会悄悄重新生成或换字。补采与重生成后，应制作新字体并重新检查；旧评价不自动证明新结果合格。

\section{记录与复现}
工作区的 \code{generation} 目录分别保存导入、档案、结果与诊断。v6.1 的结果含 \code{protocol.json}、逐字轨迹压缩文件和来源侧记录；v4 保留轨迹目录，另有实际方法及固定配置记录。开发者字体保留生成依赖、配置、版本与哈希；用户字体保留必要版本、来源摘要和警告。完整记录不表示已打包所有原始数据，独立复现仍需要相应样本与结构资源。

排版后如需查看来源缺口，在排版页打开\ui{设置 → 问题 → 字体提示}。主窗口不再单独显示跨模块工具栏；字体库直接通过同名标签页打开。

\section{字体库与生成页的分工}
生成页按准备样本、生成汉字、生成结果组织操作。运行记录默认折叠，需要查看详细路径和处理原因时展开。字体库默认只展示已经制作或导入的字体，可以用于排版、导出用户版或验证。

从字体库点击\ui{制作字体…}，也可选择历史生成结果或仅添加符号样本。制作窗口不会调用生成器；它将已有成品轨迹封装为字体。制作过程中输入锁定，成功后自动选中新字体。跨作者交接时不会沿用另一作者的符号目录。

\section{按字集与多进程（0.1.40）}
\ui{内置简繁字表…}提供大陆《通用规范汉字表》2013 年版 8105 字、GB/T 12345-1990 繁体汉字 6866 字、台湾《常用國字標準字體表》1982 年版 4808 字，可自由合并。另有按 Unihan 17.0.0 全对应构造的 8593 字派生繁体扩展集，不属于地区规范。显示来源与版本，合并去重保留共用字；一简多繁、自身对应及地区编码异体保留独立身份，不进行用户文本转换。

结构库缺字完整列出，可复制报告。只有明确勾选已查看缺口并仅生成支持部分后，才允许排除缺口；改变选择须重新确认。完整原表、实际目标、缺口、版本与哈希均随任务保存。选择字表不改变供体池，也不保证同一编码的地区笔形；本版没有香港规范完整字表。台湾源表 A02149 的编码栏冲突有明确勘误及核对依据。

\ui{导入字表…}仍支持 UTF-8 TXT 或 CSV；多列 CSV 必须有唯一的 \code{character}、汉字或字符列。去重保留首次顺序，不做简繁转换，不拆开变体序列。自定义字表不支持的字符仍阻止任务启动。编辑内置字表后改为自定义字表，原字表身份不再沿用。

\ui{从文档提取汉字…}支持 Markdown 和已有项目；\ui{从当前排版提取}读取当前未保存源码及模板配置，不改原文、设置或冻结选样。读取正文、公式、代码和图题中的汉字，不读 URL 或图片路径；当前模板只纳入活动布局引用的字段。没有模板上下文的 Markdown 采用保守文本集合，原解析诊断保留。提取结果仍可编辑。

\ui{计算方式}提供自动、串行、二进程和四进程。自动模式少于 128 字使用串行，大任务同时受可用内存和 CPU 限制，上限四进程；v4 始终串行。按字分派，每个进程复用完整作者档案，不缩小供体池，不省略语义、碰撞或来源检查。

v6.1/v6.2 任务创建独立 SQLite 一致性备份，包含已提交 WAL 内容；逻辑身份只在此不可变资源视图内复用。新任务重新读取，不按路径或修改时间跨任务缓存。\code{target-manifest.json} 保存实际字符与字表身份；\code{protocol.json} 保存规则、资源、计算方式、失败、取消及待完成项。

工作进程不可用时，记录原因并串行补齐未提交字符；已有完成字不重算。只承认逐字完整提交与哈希一致的结果，残留暂存文件不等于完成。取消后只保留已提交结果；尚未提供自动断点续跑。制作字体仍另存新包，旧候选、旧文档与冻结选样不自动迁移。

\part{书写采集}
\chapter{采集与资源交接}\label{ch:capture}
\section{书写者与任务管理（0.1.37）}
先选择或新建书写者，新建只填写显示名称，系统自动生成稳定唯一身份；旧 user\_001 原样保留，修改显示名称不改样本归属。汉字与扩展共用作者列表，启动不自动打开任务。导入 TXT/CSV 先集中确认作者、字表、份数与目录；取消不替换当前任务。

\ui{暂停并关闭}保留任务进度；\ui{取消任务}只停止剩余采集，已有样本不删除；在\ui{任务列表}恢复后可继续。列表统一展示汉字、扩展、已完成、已取消及损坏任务，可搜索，支持修改任务名称、归档、复制未录部分到另一个作者。复制不继承原作者的复用确认。零样本任务可经确认更换作者，已有样本不直接改归属。

两类采集均可查看和重录已保存份次，保留原始校验备份，汉字对应旧 NPY/TXT 也从活动数据中退役。已有字体、生成结果和冻结文档不自动改写。未保存笔迹提供保存、放弃或返回；多人窗口发生进度冲突时拒绝覆盖并保留笔迹。Windows 数据目录过深时提示改选较短目录。以下历史操作中填写 ID 的步骤，在当前统一应用中使用书写者选择。
\entry{统一应用左侧\ui{采集}导航，内含两个子页：汉字样本、标点与扩展。}

\section{汉字采集}
填写作者 ID，创建或打开采集任务，使用数位板书写，核对笔迹后保存。汉字单字任务保留标准笔数和笔序检查。画布使用数位笔输入；鼠标可操作按钮和表单，不能代替数位笔录制。应用不再显示整行采集页，已有整行资源仍保留。

完成后切换生成页，填写作者 ID、选择汉字样本目录，依次导入、建档、生成。导入会检查作者身份，不能仅修改界面的作者名来混用其他作者样本。

\section{标点与扩展}
目标字符使用大字预览，正体与斜体分别显示。误录且已经保存时，点击\ui{重录已保存…}，查看旧笔迹并选择对应份次；可以批量选择同字符、同写法的全部已保存记录。确认后所选份次退回待采，原始记录保留备份，正确样本不变。重新保存产生新样本，不会把旧错误记录再次加入字体。\ui{清空当前笔迹}只清除未保存的画布内容。已完成任务可通过\ui{继续／打开任务}重新打开纠错。

历史备份在采集目录的\code{symbol\_collection/revisions}，不参与字体制作。已有字体和冻结文档保持原样；修改采集样本后需要重新制作字体才能采用新笔迹。

类别未勾选时显示空心灰框，勾选后显示系统强调色底和白色对勾，便于在深色与浅色主题下识别。

符号目录按英文写法、数字、希腊字母及变体、括号书名号与定界符、标点连接符、单位量值和数学符号分组。类别行显示代表字符，展开可逐项选择。十三对明确列出的全角、半角同形符号各只录一次。可通过指定共用供体预览来源，并由用户逐项确认。

制作字体时自动为共用标点建立各自 Unicode 字符，笔迹不拉伸，全角形式保留全角前进宽度。旧样本可复用，旧字体与冻结文档不会自动改写；使用更新后的字体需在新版重新制作。完整字体同时选择汉字生成结果和符号样本，不要用仅符号的验收字体替代完整字体。

打开\ui{标点与扩展}，点击\ui{创建任务}。在弹窗中填写任务名称、作者 ID，勾选\ui{英文正体}、\ui{英文斜体}、数字、运算符、希腊字母、希腊变体、括号或标点等类别；设置每字符份数，默认五份，可设一至一百份。下方显示字符/写法数与总样本数；未选择类别时不能创建。英文正体和斜体可同时加入同一任务，同一字母的两种写法分别保存，继续任务后保持分类。取消配置不会替换当前任务。导入 TXT/CSV 时保留字表份数，TXT 与 CSV 空份数默认五份。画布大小可以调整，存在未保存笔迹时先保存或清空。符号采集结果通过字体库加入同一作者的字体，无需经过汉字生成。

正文优先选择正体，公式变量优先选择斜体；没有相应类别会提示，应用不会自动倾斜笔迹来伪造实录斜体。非标准笔数实录可以用于排版，不代表它具有可靠的标准笔画映射或部件生成资格。

\section{备份与设备验证}
保留原始采集 JSON 和任务记录，跨机器转移时确认作者 ID 与目录完整。生成结果和字体发布包分别保存，不覆盖原始样本。真实数位板压力、设备驱动及写字机纸面效果需要实际硬件验证；模拟采集检查不能代替这些验收。Windows 统一应用的构建与验收结果以对应版本记录为准。

\section{自由组合与确认复用（0.1.10）}
创建任务不提供快捷预设。各类字符默认不选，可整类勾选或展开逐字选择；英文大小写分别显示，斜体仅选需要补录的字母。默认每字五份，已有活动样本按同作者、字符和写法补差额；不会重建旧任务或清除进度。

点击\ui{查看笔迹并配置复用…}，查看每份已有英文正体笔迹，逐项确认斜体复用或相似希腊大写字符复用。未选录并不表示复用。未录来源须先采集；取消不保存，切换作者清除未提交关系。也可仅保存复用配置。制作字体须选择含 tasks 和 samples 的完整采集目录或 ZIP；来源缺失或重录时需重新确认。独立录制的目标字形优先，原始笔迹不变形，希腊字符保留自身身份。已有字体与冻结文档保持原样。

\section{公式数学轴信息（0.1.11）}
从原始样本制作字体时，数学轴参考信息保存在受哈希保护的核心来源摘要中，开发者转用户版和排版项目保存重开均保留。未审阅校准提示仍显示。旧字体如果缺少数学轴，需选择原始符号样本重新制作；无需重新录制。应用不会凭空填入默认轴值，也不会静默更新旧字体或冻结项目。

\appendix
\part*{附录}
\addcontentsline{toc}{part}{附录}
\chapter{常见问题与支持范围}\label{ch:trouble}
\section{遇到问题先看哪里}
\begin{longtable}{@{}P{39mm}P{115mm}@{}}
\toprule
现象 & 处理方法\\\midrule\endhead
左侧输入没有文字，撤销落到 Finder & 先确认运行的是焦点修复版，保存并退出旧副本；重新打开后确认菜单栏属于 InkStruct，再点源码区。若仍异常，记录应用路径、版本、输入法及操作步骤。\\
左侧有字，右侧没有对应笔迹 & 等待最新预览；点击查看问题，检查字形资源、语法和溢出。系统字体显示的输入文字不是手写资源。\\
预览看起来旧，导出按钮禁用 & 当前内容可能仍在重排或已有阻断错误；以最新状态与诊断为准，不使用旧页作为完成结果。\\
设置某张表格 & 打开元素中的表格设置，在应用范围选择目标表格；尚无表格时可先配置文档默认。\\
数字变大后压到框线 & 区分数字高度、最小/固定行高与留白，允许行高增长或调整字号；不要只改屏幕缩放。\\
公式显示未知命令 & 检查受控语法、花括号及环境闭合；不能直接粘贴完整 LaTeX 文档或宏定义。\\
录制分数线或根号失败 & 检查数学减号、根号的实际样本和字形结构，必要时显式切换几何构件。\\
公式太宽或跨行行组太高 & 拆成多个公式/表格，或调整尺寸和模板正文区；这些整体结构不能内部自动分页。\\
模板已有效，正文仍不能导出 & 几何合法不等于内容适配。检查绑定、缺字、必填字段和固定区溢出。\\
模板改名后绑定失效 & 使用联动重命名；若已经直接改了 JSON，恢复引用或重新绑定，不要忽略失效诊断。\\
有 PNG/JPEG 时机器导出禁用 & 正文位图只支持参考页面。用确认过的受控单线 SVG 才能进入机器输出。\\
另一个人修改后无法保存 & 另存自己的副本，再核对并合并需要的内容。不要把重试覆盖当作协作方案。\\
只有未标记字，是否都合格 & 不能自动判断；未标记只是没有评分。若只挑坏字标记，应在交接中注明检查策略。\\
\bottomrule
\end{longtable}

\section{本版支持范围速览}
\par\noindent\begin{tabularx}{\linewidth}{@{}P{34mm}Y@{}}
\toprule
类别 & 当前范围\\\midrule
文本 & 正文、六级标题、列表、引用、任务、链接文字、代码。\\
表格 & 管道表、受控 HTML、跨行跨列、逐列对齐、重复表头；不内部拆分跨行连接组。\\
数学 & 常用符号、上下标、分式、根式、伸缩括号、上下限、矩阵、受控多行与编号引用；不执行任意 TeX。\\
图形 & 位图参考页面；经校验的单线 SVG 可写。\\
模板 & 首页/续页、单正文流、固定区与字段、底图、平移校准；不支持多栏及任意页型。\\
资源与保存 & 已有轨迹、录制符号、稳定选样、资产评价、内嵌项目与会话恢复。\\
输出 & 参考 SVG/PNG、逐页机器 SVG、评价 CSV；不是设备控制台。\\
\bottomrule
\end{tabularx}

\section{反馈问题时附带什么}
先保存可重开的项目副本，再说明复现步骤、期待结果、实际结果以及涉及的页码。界面问题注明应用包名称、系统和输入法；字形问题注明作者、资产版本及具体字符；模板问题附模板 JSON 或适配报告。屏幕截图有助于说明现象，项目与报告有助于复现。

需要检查纸面偏移时，另外记录纸张尺寸、模板模式、输出缩放、笔具与设备设置。手册、软件检查和屏幕预览都不能单独证明某台写字机的纸面效果。

\chapter{Agent 协作与可复用交接}\label{ch:agent}
本附录供希望让 Agent 设计模板、检查缺口或批量导出的用户阅读。普通编辑不需要使用命令行。Agent 应从已保存项目和明确模板开始，使用应用的同一套校验规则。

\section{给 Agent 的任务信息}
可以提供以下要求：纸张实际宽高；预印纸还是空白纸；首页和续页可写区域；哪些字段固定、必填或每页重复；字号、表宽和行高要求；现有项目和资源；期望交付模板、项目、预览还是机器 SVG。

可使用这样的任务描述：
\begin{quote}
为这份已保存项目设计一个通用记录单模板。纸张 A4，预印模式；首页保留两个固定填写区，剩余内容自动续页。请先列出字段与源码块的映射，再完成绑定和实际内容试排，交付模板 JSON、新项目、预览及适配报告。保留原项目和原始资源。
\end{quote}
不要把“参考某张学校表单”自动理解为软件必须内置那所学校的栏目。模板可以依据参考材料设计，但应保存成独立可复用配置。

\section{命令入口约定}
以下示例以已安装环境的 \code{inkstruct} 命令为例。仓库用户从自己的虚拟环境调用该程序；macOS 独立包可使用下面的前缀替代命令名：
\par\noindent\begin{minipage}{\linewidth}
\begin{Source}
"/实际位置/InkStructEditor.app/Contents/MacOS/InkStructTool" --inkstruct-cli
\end{Source}
\end{minipage}\par\medskip
示例中的文件名代表用户实际文件；输出文件或目录应尚不存在。命令运行成功后仍要查看报告，不只依据退出状态判断机器可写。

\section{模板设计与内容绑定的顺序}
\par\noindent\begin{minipage}{\linewidth}
\begin{Source}
inkstruct document template new --output layout.json
inkstruct document template validate layout.json
inkstruct document template preview layout.json --output layout-preview
inkstruct document template inspect-slots layout.json
inkstruct document template inspect-content report.inkdoc --template layout.json
\end{Source}
\end{minipage}\par\medskip
新建模板后，Agent 修改 JSON 副本并再次校验。几何预览只展示区域和辅助文字，不能当作机器 SVG。使用语义字段时，从 \code{inspect-content} 输出提取 \code{slot_request} 草稿，填入实际字段与本次返回的源码块编号，另存为请求文件。

\par\noindent\begin{minipage}{\linewidth}
\begin{Source}
inkstruct document template bind report.inkdoc --template layout.json --bindings request.json --output bound.inkdoc
inkstruct document template validate-document bound.inkdoc
inkstruct document render bound.inkdoc --plotter --optimize-paths --output machine-output
\end{Source}
\end{minipage}\par\medskip
请求中的项目哈希与节点编号来自同一次检查，不能猜测。绑定请求会\textbf{全部替换}已有绑定，要保留的项也必须列出；未列出的块返回正文流。项目在检查后发生变化，应重新读取并生成请求。

没有语义字段的模板可使用返回的普通 \code{binding_request}。两种请求不能随意拼接字段。存在错误时先修正，再保存新项目；不可书写的位图即使内容适配通过，也仍不能机器导出。

\section{检查补采需求}
\par\noindent\begin{minipage}{\linewidth}
\begin{Source}
inkstruct document resources report.inkdoc --capture-csv missing-symbols.csv
\end{Source}
\end{minipage}\par\medskip
原生项目继承内嵌资源与公式构件来源。报告分开列出可采集符号、汉字和其他不支持缺口，符号 CSV 每字 5 次。这个命令只读检查资源，不修改项目，也不验证纸面几何或根号能否识别。补充资源后需重新试排。

\section{交付与复核}
Agent 应交付修改后的模板 JSON、新 \code{.inkdoc}、当前预览及适配报告；要求机器输出时另附逐页 SVG 和清单。说明仍缺的资源、校准状态及没有完成的验证。输入项目、原始样本和参考 PDF 保留不变。


\end{document}
